% =====================================================================
\section{Background and related work}
\label{sec:related}
% =====================================================================

\paragraph{Multiple adaptation loops}
Self-adaptive systems are commonly built from feedback loops of the
monitor--analyse--plan--execute kind~\cite{kephart2003}, and a system with
several adaptation concerns may run one loop per concern. Weyns
et~al.\ catalogue patterns for organising such loops, from fully
decentralised to hierarchical control~\cite{weyns2013}. That uncoordinated
loops can interfere is well established: independent data-centre power
managers fight one another~\cite{nopowerstruggles}, independent
per-service autoscalers oscillate where a coordinated one does
not~\cite{chamulteon}, and Kephart et~al.\ coordinated a performance
manager and a power manager through one shared signal~\cite{kephart2007}.
In process control, Bristol's relative gain array measures when
decentralised single-loop control suffices for a multivariable
plant~\cite{bristol1966}. Our width result (Section~\ref{sec:width}) is a
discrete analogue of that question for one gateway: we do not ask whether
loops can interfere, which they can, but how much a joint search recovers
when every loop is rate-limited, and we measure it with the loops' other
ingredients held fixed.

\paragraph{Serving, caching and routing for LLMs}
Inference engines such as vLLM~\cite{vllm} and Splitwise~\cite{splitwise}
manage GPU memory and execution phases inside a model server. Semantic
caches return a stored answer for a sufficiently similar
prompt~\cite{gptcache}, and their timing has been shown to leak other
users' prompts~\cite{earlybird}. Routers and cascades trade answer quality
for price per query~\cite{frugalgpt,routellm}. \polyforge{} sits above all
of these: it treats model servers as external endpoints with a price and a
latency, and decides per tenant and per interval how many gateway replicas
to run, how much cache to grant and which tier to call, under a per-tenant
dollar budget.

\paragraph{Joint configuration of serving knobs}
Several systems choose two of these knobs together. Clipper combined a
prediction cache with adaptive model selection~\cite{clipper}. INFaaS
selects model variants together with VM- and model-level
autoscaling~\cite{infaas}; IPA solves an integer program over model
variants, replicas and batch sizes for inference pipelines on
Kubernetes~\cite{ipa}; Proteus scales accuracy, choosing variants and
placement under load~\cite{proteus}. For LLMs, Cascadia co-optimises the
deployment of a model cascade with its routing policy~\cite{cascadia},
Cheng and Nguyen jointly choose models, GPU provisioning, parallelism and
routing under latency, accuracy and budget constraints~\cite{chengnguyen},
and Zhu et~al.\ design a cache jointly with a model selector~\cite{zhu2023}.
Each reports a gain over separately configured or fixed alternatives. None
of them, to our knowledge, separates coordination from the rest of the
joint design by an ablation that keeps the objective, forecaster, budgets
and move limits while searching the knobs separately, and none controls the
replica, semantic-cache and tier triple of a multi-tenant gateway. Our
contribution is not a better joint controller; it is a measurement of how
much the jointness of one is worth, and of what else its apparent wins
were made of.

\paragraph{Evaluation methodology}
Autoscaler comparisons need a level playing field~\cite{chameleon}.
Recent work makes the same point in our setting: a calibrated rule-based
scaler beat six deep reinforcement learning methods on
cost~\cite{zhang2026drl}, and a pre-registered re-evaluation of learned
orchestrators traced apparent robustness to comparator collapse and
evaluation choices~\cite{khilji2026}. In machine learning, re-evaluations
under equal tuning budgets erased claimed gains in language
modelling~\cite{melis2018}, generative adversarial networks~\cite{lucic2018}
and deep reinforcement learning~\cite{henderson2018}, and Hoefler and Belli
set out reporting rules for performance measurements~\cite{hoefler2015}.
Pre-registration, committing hypotheses, design and analysis before data
exist, separates confirmatory from exploratory claims~\cite{nosek2018}.
We apply these practices to our own controller: every campaign reported
here has a protocol committed before its first run, and every correction
of a comparator is a new protocol beside the old record, never an edit of
it.
